\scanpage{59}
\mspage{- 26 -}

In $\underline{S}$, roots exist, and hence this natural transformation
induces a unique map:
\[
\varphi_{G,D}:\mathrm{Rr}\operatorname{Hom}_{\underline{A}}(G,D(-))
\longrightarrow\operatorname{Hom}_{\underline{A}}(G,\mathrm{Rr}\,D).
\]
We are interested in the case where $\underline{D}$ is a directed system.
(See \S 4, Ex. iii).) Let $D(i)=A_i$, $i\in\underline{D}$. Then
\[
\varphi_{G,D}:\operatorname{dir\,lim}\operatorname{Hom}_{\underline{A}}(G.A_i)
\longrightarrow\operatorname{Hom}_{\underline{A}}(G,\operatorname{dir.\,lim}A_i).
\]

\uline{5.2.} \uline{Definition:} $G\in\underline{A}$ is called \uline{small} if for every directed
system $D:\underline{D}\longrightarrow\underline{A}$, the above morphism $\varphi_{G,D}$ is a monomorphism.
If there is a generating family (\S 1,vii) of small objects, then
we say that $\underline{A}$ \uline{has enough} small objects. \uline{Cosmall} objects are
defined dually.

\uline{5.3.} \uline{Examples:} In the category of sets or topological spaces,
the set consisting of a single point is a small generator. In
the category of abelian groups, $Z$ is a small generator. In
the category of torsion abelian groups, $\{Z_p\}$ is a generating
family of small objects. In the category of modules over $R$,
$R$ is a small generator. Dual statements hold in the dual
categories. For a characterization of small objects in special
categories, see 6.4. below.
\mspage{58}
